\section{\methodname}
\label{sec:switch}

\subsection{\tensorswitchplus}
\label{sec:tensorswitch}
A limitation of the above \blazeplus construction is that its query complexity
scales linearly with $t$.
This overhead is unavoidable when using a point oracle
$\Oracle_{\mathsf{P}}(\bm{C})$, since each selected row must be queried entrywise.

To reduce this cost, we replace the point oracle for $\bm{C}$ with an
MLE oracle for a $\mu$-wise interleaved code,
$\Oracle_{\mathsf{MLE}}[\Code_2^{t}](\bm{C})$.
Here, each of the $\mu$ rows of $\bm{C}$ is encoded using a linear code
$\Code_2:\Field^{t}\to\Field^{\mu_2}$,
and all accesses to $\bm{C}$ are performed via MLE queries to
$\Oracle_{\mathsf{MLE}}[\Code_2^{t}](\bm{C})$.
Observe that Algorithm~\ref{alg:interleave} only queries $\bm{C}$ through MLE
evaluations on selected rows.
Therefore, by using the oracle
$\Oracle_{\mathsf{MLE}}[\Code_2^{t}](\bm{C}^\top)$,
these queries can be answered succinctly.

By Fact~\ref{fact:tensorAsInterleave},
the MLE oracle for $\bm{C}=\bm{G}\bm{M}$ corresponds to the tensor code
$\Code \otimes \Code_2$,
and thus we obtain an MLE oracle for a tensor code, denote by
\tensorswitchplus.

The construction of \tensorswitchplus modifies
Algorithm~\ref{alg:interleave} as follows:
\begin{itemize}
    \item The input oracle is replaced by the point oracle
    $\Oracle_{\mathsf{P}}(\bm{C}\bm{G}_2^\top)$.
    \item In Line~\ref{line:interleaveSendv}, the prover sends an MLE oracle
    $\Oracle_{\mathsf{MLE}}[\Code_2^{k/t}](\bm{v})$.
    \item In Lines~\ref{line:interleaveQueryC} and~\ref{line:interleaveRowMLE},
    instead of querying individual entries of $\bm{C}_{\mathcal{I}}$,
    the verifier queries $\Oracle_{\mathsf{P}}(\bm{C}\bm{G}_2^\top)$
    for the Merkle roots of the selected rows.
    \item The verifier combines the Merkle root of
    $\Oracle_{\mathsf{MLE}}[\Code_2^{k/t}](\bm{v})$
    with that of the selected rows to obtain the root of
    $\Oracle_{\mathsf{MLE}}[\Code_2^{k/t+\ell}]
    (\bm{v}\,||\,\mathsf{vec}(\bm{C}_{\mathcal{I}}))$.
    \item Finally, the verifier queries to the MLE oracle
    $\Oracle_{\mathsf{MLE}}[\Code_2^{k/t+\ell}]
    (\bm{v}\,||\,\mathsf{vec}(\bm{C}_{\mathcal{I}}))$ to verify the matrix-function relation:
    \[
        v' + z\hat{v}' + z^2\,\bm{\xi}^{\top}\bm{\eta}\bm{C}_{\mathcal{I}}\tilde{\bm{r}}
        \EqOrNot
        \langle
            \bm{l} + z\hat{\bm{l}} + z^2\,\bm{\xi}^{\top}\bm{\eta}\bm{G},
            \bm{v}
        \rangle,
    \]
    where the MLE query is switched to an MLE-friendly code
    $\Code':\Field^{k+\ell t}\to\Field^{\mu'}$.
\end{itemize}

This process is illustrated in Fig.~\ref{fig:enc}.
For brevity, we omit the full pseudocode of \tensorswitchplus;
it is incorporated as a subroutine of the \methodname protocol
in the next section.

\begin{theorem}
\label{theorem:tensorswitch}
The $\tensorswitchplus$ MLE oracle constructed above is computationally
and knowledge sound, with total error bounded by
\[
    \frac{2\log_2 k}{|\Field|-2} + \exp(-2\iota^2 \cdot \ell_1)  +  \frac{\log_2 \ell_2 +1}{|\Field|-2} + (1-(\epsilon_1 - \iota)\cdot \epsilon_2)^{\ell_2} +
    + \kappa[\Code']
    + \kappa_{\mathsf{MFO}}
\]
where $\kappa[\Code^{t}]$, $\kappa[\Code_2^{k/t+\ell}]$, $\kappa_{\mathsf{MFO}}$, and $\kappa[\Code']$
denote the computational/knowledge soundness errors of Algorithm~\ref{alg:interleave} for $\Code^t$ or $\Code_2^{k/t+\ell}$,
the matrix-function oracle, and the code-switching procedure
in Algorithm~\ref{alg:basicMLEOracle}, respectively.
\end{theorem}



\begin{proof}
The claim follows from
Theorem~\ref{theorem:interleavedIOPP},
Theorem~\ref{theorem:iop},
and Lemma~\ref{lemma:codeswitch}.
\end{proof}




\subsection{Filtering and Combining Point Oracles}

\subsection{\tensorswitchplus}
\label{sec:tensorswitch}
$\Code_2: \Field^t \mapsto \Field^{\mu_2}$


\begin{theorem}
\label{theorem:tensorSwitch}
Given a linear code $\Code: \Field^k \mapsto \Field^\mu$ with generator matrix $\bm{G} \in \Field^{\mu\times k}$ and relative distance $\delta$, its tensor code $\Code^{\kronecker 2}: \Field^{k^2} \mapsto \Field^{n^2}$,
and another linear code $\Code': k \mapsto \mu'$ with relative distance $\delta'< \delta$ that admits an MLE IOPP for any proximity parameter $\epsilon' \le \delta/3$.
Given a proximity parameter $\epsilon \le \delta/3$,
there exists an MLE IOPP for $\Code^{\kronecker 2}$.
\end{theorem}

\medskip
\noindent
{\bf Query complexity.}
The query complexity of \tensorswitchplus contains three part:
\begin{itemize}
    \item The query of $\ell$ rows for interleaved-code IOPP, requriing $\ell$ Merkle path of height $\log_2 \mu$ from the row root to the tensor codeword Merkle root.
    \item The query of $\ell'$ positions for code switching to MLE oracle of $\Code'$, where the Merkle tree height is
    $\log_2 [(k/t+\ell)\mu_2]$.
    \item And the query complexity of $\Code'$ for a message of lengh ${k/t+\ell}$.
\end{itemize}

the query of $\ell'$ positions for code switching to MLE oracle of $\Code'$,
and the query complextiy of $\Oracle_\mathsf{MLE}[\Code']$.

Let us set both $\Code$ and $\Code_2$ to be interleaved RAA code with rate-$1/4$ and relative distacne $\delta = 0.19$,
as per Lemma~\ref{corollary:interleaveRadius},
any interleaved code of this RAA code has list-decoding radius $\delta_\mathsf{LD} = 0.19/2$,
hence $\ell = \ell' = 7\lambda$
We choose $\Oracle_\mathsf{MLE}[\Code']$ to be an low Query complexity MLPC,
\eg, WHIR~\cite{eurocrypt/ArnonCFY25} with $O(\log\log n \cdot \lambda)$ query complexity.
So the overall 

\input{alg/tensor_switch}


\subsection{\methodname}
\label{sec:dualswitch}
\input{alg/dualswitch}

\begin{proof}
Let $\bm{C}:=\bm{C}_1-\bm{C}_2\in\Code^{\kronecker 2}$.
Since $\bm{C}_1\neq\bm{C}_2$, we have $\bm{C}\neq\bm{0}$.
Therefore, $\bm{C}$ contains at least one nonzero column.
Because each column of $\bm{C}$ is a codeword of $\Code$,
it has Hamming weight at least $\delta\mu$.

As a consequence, $\bm{C}$ has at least $\delta\mu$ nonzero rows.
Each such nonzero row is itself a codeword of $\Code$
and thus has Hamming weight at least $\delta\mu$.
Hence:
\[
\left|
\left\{
i\in[\mu]
\;\middle|\;
\Delta_H(\bm{c}_{i*},\bm{0}) \ge \delta\mu
\right\}
\right|
\;\ge\;
\delta\mu,
\]
which implies
$\delta^{\kronecker 2}_{\delta}(\bm{C}_1,\bm{C}_2)\ge\delta$.
\end{proof}




\begin{proof}
\emph{Completeness.}
From the completeness of inner product proof, we know the inner-product proofs passed.

Moreover we know:
\begin{align*}
& \bm{C} = \bm{G}_1\bm{M}\bm{G}_2^\top
~\implies~
 \bm{\eta}_L \bm{C} \bm{\eta}_R = \bm{\eta}_L\bm{G}_1\bm{M}\bm{G}_2^\top\bm{\eta}_R
\\&
~\implies~
 \langle \bm{\xi}, \mathsf{\diag}(\bm{\eta}_L\bm{\eta}\bm{C}\bm{\eta}_R^\top) \rangle
=\langle \bm{\xi}, \mathsf{\diag}(\bm{U}\bm{G}^\top\bm{\eta}_R^\top) \rangle
\end{align*}
And hence we got completeness.

\medskip \noindent
\emph{Soundness.}
We define the following six state, where the meaning of the variables are those in Algorithm~\ref{alg:tensorSwitch}.
We prove soundness from round-by-round soundness between these states.

{\allowdisplaybreaks
\begin{align*}
\mathsf{State}_0&
(
\bm{C}, \hat{v}, \hat{\bm{l}}, \hat{\bm{r}}, v, \bm{l}, \bm{r}: \bm{M}
)
\\
=\{&
~v = \bm{l}^\top \bm{M} \bm{r}
~\wedge~
\hat{v} = \hat{\bm{l}}^\top\bm{M} \hat{\bm{r}}
~\wedge~
\delta_{\epsilon_2}^{\kronecker 2}(\bm{C}, \bm{G}\bm{M}\bm{G}^\top) \le \epsilon_1
\},
\\
%%%%%%%%%%%%%%%%%%%%%%%%%%%%%%
\mathsf{State}_1&
(
\bm{C}, \hat{v}, \hat{\bm{l}}, \hat{\bm{r}}, v, \bm{l}, \bm{r}, \tilde{\bm{r}},\hat{v}', v': \bm{M}
: \bm{M}
)
\\
=\{&
~\text{inner-product proofs pass}
\\
&
~\wedge~
v' = \bm{l}^\top \bm{M} \tilde{\bm{r}}
~\wedge~
\hat{v}' = \hat{\bm{l}}^\top\bm{M} \tilde{\bm{r}}
~\wedge~
\delta_{\epsilon_2}^{\kronecker 2}(\bm{C}, \bm{G}\bm{M}\bm{G}^\top) \le \epsilon_1
\},
\\
%%%%%%%%%%%%%%%%%%%%%%%%%%%%
\mathsf{State}_2&
(
\bm{C}, \hat{v}, \hat{\bm{l}}, \hat{\bm{r}}, v, \bm{l}, \bm{r}, \tilde{\bm{r}},\hat{v}', v', \bm{v}: \bm{M}
)
\\
=&\{
~\text{inner-product proofs passed}
\\
&
~\wedge~
\bm{v} = \bm{M}\tilde{\bm{r}}
~\wedge~
v' = \bm{l}^\top \bm{v}
~\wedge~
\hat{v}' = \hat{\bm{l}}^\top\bm{v} 
~\wedge~
\delta_{\epsilon_2}^{\kronecker 2}(\bm{C}, \bm{G}\bm{M}\bm{G}^\top) \le \epsilon_1
\},
\\
%%%%%%%%%%%%%%%%%%%%%%%%%%%%
\mathsf{State}_3&
(
\bm{C}, \hat{v}, \hat{\bm{l}}, \hat{\bm{r}}, v, \bm{l}, \bm{r}, \tilde{\bm{r}},\hat{v}', v', \bm{v}, \bm{\eta}, \bm{U}
: \bm{M}
)
\\
=&\{
~\text{inner-product proofs passed} ~\wedge~
v' = \bm{l}^\top \bm{v}
~\wedge~
\hat{v}' = \hat{\bm{l}}^\top\bm{v} 
\\
&
~\wedge~
\bm{\eta}_L\bm{G}\bm{v} = \bm{U}\bm{r}'
~\wedge~
\bm{\eta}_L\bm{G}\hat{\bm{v}} = \bm{U}\hat{\bm{r}}
~\wedge~
\delta_{\epsilon_2}^{\kronecker 2}(\bm{\eta}\bm{C}, \bm{U}\bm{G}^\top) \le \epsilon_1 - \iota
\},
\\
%%%%%%%%%%%%%%%%%%%%%%%%%%%%
\mathsf{State}_4&
(
\bm{C}, \hat{v}, \hat{\bm{l}}, \hat{\bm{r}}, v, \bm{l}, \bm{r}, \tilde{\bm{r}},\hat{v}', v', \bm{v}, \bm{\eta}, \bm{U}, \bm{\eta}_L, \bm{\eta}_R
: \bm{M}
)
\\
&
~\text{inner-product proofs passed}
~\wedge~
v' = \bm{l}^\top \bm{v}
~\wedge~
\hat{v}' = \hat{\bm{l}}^\top\bm{v} 
\\
&
~\wedge~
\bm{\eta}_L\bm{G}\bm{v} = \bm{U}\tilde{\bm{r}}
~\wedge~
\mathsf{\diag}(\bm{\eta}_L\bm{\eta}\bm{C}\bm{\eta}_R^\top)
=\mathsf{\diag}(\bm{U}\bm{G}^\top\bm{\eta}_R^\top)
\},
\\
%%%%%%%%%%%%%%%%%%%%%%%%%%%%
\mathsf{State}_5&
(
\bm{C}, \hat{v}, \hat{\bm{l}}, \hat{\bm{r}}, v, \bm{l}, \bm{r}, \tilde{\bm{r}},\hat{v}', v', \bm{v}, \bm{\eta}, \bm{U}, \bm{\eta}_L, \bm{\eta}_R, \bm{\xi}, z
: \bm{M}
)
\\
=&\{
~\text{inner-product proofs passed}
~\wedge~
(v'+z\hat{v}') = (\bm{l}^\top + z\hat{\bm{l}}^\top) \bm{v}
\\
&
~\wedge~
\bm{\xi}^\top\bm{\eta}_L\bm{G}\bm{v} = \bm{\xi}^\top \bm{U}\tilde{\bm{r}}
~\wedge~
\langle \bm{\xi}, \mathsf{\diag}(\bm{\eta}_L\bm{\eta}\bm{C}\bm{\eta}_R^\top) \rangle
=\langle \bm{\xi}, \mathsf{\diag}(\bm{U}\bm{G}^\top\bm{\eta}_R^\top) \rangle
\},
\end{align*}
}



\smallskip
\noindent
{\sf $\mathsf{State0} \rightarrow \mathsf{State1}$.}

If 
\(
v \ne \bm{l}^\top \bm{M} \bm{r} \vee \hat{v} \ne \hat{\bm{l}}^\top \bm{M} \hat{\bm{r}}
\),
then due to the soundness of the sumcheck protocol,
the following event happens with probability at most $\frac{2\log_2 k}{|\Field|-2}$
\begin{align*}
&
    \Pr \left[
        \text{inner-product proofs pass} \wedge 
        v' = \bm{l}^\top \bm{M} \tilde{\bm{r}} \wedge
        \hat{v}' = \hat{\bm{l}}^\top \bm{M} \tilde{\hat{\bm{r}}}
        \left|
        v \ne \bm{l}^\top \bm{M} \bm{r} \vee \hat{v} \ne \hat{\bm{l}}^\top \bm{M} \hat{\bm{r}}
        \right.
    \right]
\\
&
    \le \frac{2\log_2 k_2}{|\Field|-2}.
\end{align*}

Hence,
\[
    \Pr_{\bm{r}}\left[ \mathsf{State1} \mid \neg \mathsf{State0} \right] \le  \frac{2\log_2 k_2}{|\Field|-2}.
\]

\smallskip
\noindent
{\sf $\mathsf{State1} \rightarrow \mathsf{State2}$.}
We know:
\[
    v' \ne \bm{l}^\top \bm{M} \tilde{\bm{r}} \wedge \bm{v} = \bm{M}\tilde{\bm{r}}
    \implies
    v' \ne \bm{l}^\top \bm{v},
\]
\[
    \hat{v}' \ne \hat{\bm{l}}^\top \bm{M} \tilde{\bm{r}} \wedge \bm{v} = \bm{M}\tilde{\bm{r}}
    \implies
    \hat{v}' \ne \hat{\bm{l}}^\top \bm{v}.
\]

Hence,
\[
    \Pr_{\bm{\eta}}\left[ \mathsf{State2} \mid \neg \mathsf{State1} \right] = 0.
\]

\noindent
{$\mathsf{State2} \rightarrow \mathsf{State3}$.}

We assume $\epsilon_1 \le \delta_\mathsf{LD1}, \epsilon_2  \le \delta_\mathsf{LD2}$.

If $\bm{v} \ne \bm{M}\tilde{\bm{r}}$, then $\bm{G}\bm{v} \ne \bm{G}\bm{M}\tilde{\bm{r}}$.
For two codeword, there distance is no less than $\delta \cdot \mu$.
Then for random $\ell$ rows, the probability that they are all zero is bounded by:
\[
    \Pr_{\bm{\eta}} \left[ \bm{\eta}_L\bm{G}\bm{v} = \bm{\eta}\bm{G}\bm{M}\tilde{\bm{r}} \mid \bm{v} \ne \bm{M}\tilde{\bm{r}} \right]
    \le (1-\delta)^\ell
    \le (1-\epsilon_2)^\ell
\]


If $\delta_{\epsilon_2}^{\kronecker 2}(\hat{\bm{u}}_r, \bm{G}\bm{M}\bm{G}^\top\hat{\bm{r}}) > \epsilon_1$,
and we call the $i$th row of $\bm{C}$ a \emph{bad} row
if its Hamming distance with the $i$th row of $\bm{G}^\top \bm{M} \bm{G}$ exceeds $\epsilon_2 \mu$.
As per the Chernoff-Hoeffding inequality,
for randomly selected $\ell$ rows via $\bm{\eta}_L$,
the proability that the number of bad rows less than $(\epsilon_2 - \iota) \cdot \ell$ is small.
Combining this with the definition of $\delta_{\epsilon_1}^{\kronecker 2}$,
we know:
\[
    \Pr_{\bm{\eta}_L} \left[
        \delta_{\epsilon_2}^{\kronecker 2}(\bm{\eta}\bm{C}, \bm{\eta}\bm{G}\bm{M}\bm{G}^\top) \le \epsilon_1 - \iota
        \mid
        \delta_{\epsilon_2}^{\kronecker 2}(\bm{C}, \bm{G}\bm{M}\bm{G}^\top) > \epsilon_1
    \right]
    \le \exp (-2 \iota^2 \cdot \ell_1).
\]

Combining the above:
\[
    \Pr_{\bm{\eta}}\left[ \mathsf{State3} \mid \neg \mathsf{State2} \right] \le  \exp(-2\iota^2 \cdot \ell_1).
\]


\noindent
{$\mathsf{State3} \rightarrow \mathsf{State4}$.}
If $\delta_{\epsilon_2}^{\kronecker 2}(\bm{\eta}_L\bm{C}, \bm{\eta}_L\bm{G}\bm{M}\bm{G}^\top) > \epsilon_1 - \iota$,
then at least for $(\epsilon_1 - \iota)\cdot \ell$ rows,
the probability of non-zero entires is at least $\epsilon_2$.

Hence, randomly selecting $\ell_2$ entries from $\bm{U}$,
the probability that all selected entries are zeros is bounded by:
\[
    (1-(\epsilon_1 - \iota)\cdot \epsilon_2)^{\ell_2}.
\]
That is:
\begin{align*}
&
    \Pr_{\bm{\eta}_R}
    \left[
        \mathsf{\diag}(\bm{\eta}_L\bm{C}\bm{\eta}_R)
        = \mathsf{\diag}(\bm{U}\bm{G}^\top\bm{\eta}_R)
        \mid
        \delta_{\epsilon_1}^{\kronecker 2}(\bm{\eta}_L\bm{C}, \bm{\eta}_L\bm{G}\bm{M}\bm{G}^\top) > \epsilon_2 - \iota
    \right]
    \\&
    \le (1-(\epsilon_1 - \iota)\cdot \epsilon_2)^{\ell_2}.
\end{align*}

Hence,
\begin{align*}
    \Pr_{\bm{\eta}}\left[ \mathsf{State5} \mid \neg \mathsf{State4} \right]
    \le   (1-(\epsilon_1 - \iota)\cdot \epsilon_2)^{\ell_2}.
\end{align*}

\noindent
{$\mathsf{State4} \rightarrow \mathsf{State5}$.}
As per the Schwartz-Zippel lemma,
\[
    \Pr_{z} \left[
    (v'+z\hat{v}') = (\bm{l}^\top + z\hat{\bm{l}}^\top) \bm{v} 
    \mid
    v' \ne \bm{l}^\top \bm{v}
    ~\vee~
    \hat{v}' \ne \hat{\bm{l}}^\top\bm{v} 
    \right]
    \le \frac{1}{|\Field|},
\]
\[
    \Pr_{\bm{\xi}} \left[
    \langle \bm{\xi}, \mathsf{\diag}(\bm{\eta}_L\bm{\eta}\bm{C}\bm{\eta}_R^\top) \rangle
=\langle \bm{\xi}, \mathsf{\diag}(\bm{U}\bm{G}^\top\bm{\eta}_R^\top) \rangle
    \mid
    \bm{\eta}_L\bm{\eta}\bm{G}\bm{v} \ne \bm{U}\tilde{\bm{r}}
    \right]
    \le \frac{\log_2 \ell}{|\Field|-2},
\]

Hence,
\[
    \Pr_{\bm{\eta}}\left[ \mathsf{State4} \mid \neg \mathsf{State3} \right] \le  \frac{\log_2 \ell_1 + 1}{|\Field|-2}.
\]



\noindent
{Summary of soundness.}
Combining the above:
\[
    \kappa^\mathsf{s} = 
    \Pr_{\bm{\eta}}\left[ \mathsf{State5} \mid \neg \mathsf{State0} \right] \le 
    \frac{2\log_2 k}{|\Field|-2} + \exp(-2\iota^2 \cdot \ell_1)  +  \frac{\log_2 \ell_2 +1}{|\Field|-2} + (1-(\epsilon_1 - \iota)\cdot \epsilon_2)^{\ell_2}.
\]

Set $\epsilon_1  = \delta_\mathsf{LD1},  \epsilon_2=\delta_\mathsf{LD2}$ and $\iota = 1/\sqrt{\log_2 \mu_1}$ to be small,
then for:
\[
    \ell_1 = \frac{2\ln 2 \cdot \lambda}{\iota^2}, \quad
    \ell_2 = \frac{\lambda}{\log_2(1-(\epsilon_1-\iota)\cdot \epsilon_2)},
\]
the protocol achieves $\lambda$-bit security.

\medskip
\noindent
{\em Knowledge Soundness.}
Given an $\Prover^*$ that produces an accepting transcript with probability $\varrho$,
use the standard tree extraction result,
an extractor $\Extractor_4$,
an extract $\bm{U}\bm{G}^\top$ that align with $\bm{\eta}_L\bm{C}$ at one position at each row.
with probability no less than:
\[
    \varrho -  \frac{\log_2 \ell +1}{|\Field|-2}
\]


\noindent
{\sf $\mathsf{State}_4 \rightarrow \mathsf{State}_3$.}
Each execution of $\Extractor_4$ extracts $\ell$ indices at which
$\bm{\eta}\bm{C}$ and $\bm{U}\bm{G}^\top$ agree.
Given such a set of sampled indices, the extractor $\Extractor_3$
can recover $\bm{U}\bm{G}^\top$ satisfying
\[
    \delta^{\kronecker 2}_{\epsilon_2}(\bm{C},\bm{U}\bm{G}^\top)
    \le \epsilon_1 - \iota,
\]
provided that the samples include at least $(1-\epsilon_2)\mu_2$ valid indices
in at least $(1-\epsilon_1+\iota)\mu_1$ rows.

Otherwise, for any invalid index set, there exist at least
$(1-\epsilon_1+\iota)\mu_1$ rows in which at most
$(1-\epsilon_2)\mu_2$ indices are acceptable.
The total number of such “bad’’ index sets is therefore bounded by
\[
    \Bigl[
        (1-\epsilon_2)\mu_2 \cdot (1-\epsilon_1+\iota)\mu_1
        + (\epsilon_1-\iota)\mu_1 \cdot \mu_2
    \Bigr]^{\ell_1}
    = \bigl[1-(\epsilon_1-\iota)\epsilon_2\bigr]^{\ell_1}.
\]
Since the total number of possible index sets is $(\mu_1\mu_2)^{\ell_1}$,
Lemma~\ref{lemma:rbrks} implies that the knowledge-soundness error
incurred in this round is at most
\[
    \bigl[1-(\epsilon_1-\iota)\epsilon_2\bigr]^{\ell_1}.
\]


\noindent
{\sf $\mathsf{State3} \rightarrow \mathsf{State2}$.}
Suppose there exists an extractor $\Extractor_3$ that,
given any transcript $\bm{\eta}_L$,
extracts $\bm U$ with probability at least
$\varrho - \kappa^{\mathsf{ks}}_3$.
In particular, with this probability,
it identifies at least $(1-\epsilon_1+\iota)\ell$ rows
for which $\bm C$ is $\epsilon_2$-close to a valid codeword.
We refer to such row indices as \emph{good indices}.

Consider one execution of $\Extractor_3$ selecting
a subset of $\ell$ indices.
Among them, at least $(1-\epsilon_1+\iota)\ell$
are good.
The number of subsets whose union contains
fewer than $\epsilon_1 \mu_1$ good indices
is bounded by:
\[
\sum_{k=0}^{(\epsilon_1-\iota)\ell}
\binom{\ell}{k}
\cdot
[(1-\epsilon_1)\mu_1]^k
\cdot
(\epsilon_1 \mu_1)^{\ell-k}.
\]

By a Chernoff bound:
\[
\sum_{k=0}^{(\epsilon_1-\iota)\ell}
\binom{\ell}{k}
\epsilon_1^k
(1-\epsilon_1)^{\ell-k}
\le
\exp\!\big(-\ell \, D(\epsilon_1-\iota \,\|\, \epsilon_1)\big)
\le
\exp(-2\iota^2 \ell),
\]
where $D(\cdot\|\cdot)$ denotes the binary
Kullback–Leibler divergence.

By Lemma~\ref{lemma:generalSpecialSound} in Appendix~\ref{appx:special},
this allows us to construct an extractor
$\Extractor_2$
that outputs $\bm M$ such that:
\[
\delta^{\kronecker 2}_{\epsilon_2}
(\bm C, \bm U \bm G^\top)
\le
\epsilon_1
\]
in expected polynomial time with knowledge soundness error
bounded by:
\[
\exp(-2\iota^2 \ell).
\]

\smallskip
\noindent
{\sf $\mathsf{State2} \rightarrow \mathsf{State0}$.}
By the knowledge soundness of the inner product proof,
the knowledge soundness error in this round is bounded by
\[
    \frac{2\log_2 k_2}{|\Field|-2}.
\]

\smallskip
\noindent
{\em Summary of knowledge soundness.}
Combining the above:
The knowledge soundness error is bounded by:
\[
    \kappa^\mathsf{ks} = 
    \frac{2\log_2 k}{|\Field|-2} + \exp(-2\iota^2 \cdot \ell_1)  +  \frac{\log_2 \ell_2 +1}{|\Field|-2} + (1-(\epsilon_1 - \iota)\cdot \epsilon_2)^{\ell_2}.
    \qedhere
\]
\end{proof}



\subsection{MLPC from \methodname}
\label{sec:mlpc}
